\begin{siglas}
  \item[DE440/441] \textit{Development Ephemeris} 440 e 441, do JPL
  \item[DFT] Transformada discreta de Fourier (\textit{Discrete Fourier Transform})
  \item[GPU] Unidade de processamento gráfico
  \item[IAU] União Astronômica Internacional
  \item[ICRF] \textit{International Celestial Reference Frame}
  \item[JPL] \textit{Jet Propulsion Laboratory} (NASA)
  \item[LRL] Vetor de Laplace--Runge--Lenz
  \item[RK4] Método de Runge--Kutta clássico de quarta ordem
  \item[TDB] Tempo Dinâmico Baricêntrico
  \item[UA] Unidade astronômica, $1\UA = 149\,597\,870{,}7\km$
  \item[WCAG] \textit{Web Content Accessibility Guidelines}
\end{siglas}

\begin{simbolos}
  \item[$G$] Constante da gravitação universal
  \item[$\mu$] Parâmetro gravitacional, $\mu = G(M+m)$, em $\mathrm{UA^3/dia^2}$
  \item[$\vr,\ \vv$] Posição e velocidade
  \item[$\vh$] Momento angular específico, $\vh = \vr\times\vv$
  \item[$\ve$] Vetor excentricidade (Laplace--Runge--Lenz normalizado)
  \item[$a,\ e,\ i$] Semieixo maior, excentricidade e inclinação
  \item[$\Omega,\ \omega$] Longitude do nodo ascendente e argumento do periélio
  \item[$M,\ E,\ \nu$] Anomalias média, excêntrica e verdadeira
  \item[$n$] Movimento médio, $n = 2\pi/T$; também o índice do harmônico
  \item[$c_n$] Coeficiente complexo de Fourier de ordem $n$
  \item[$J_n$] Função de Bessel de primeira espécie de ordem $n$
  \item[$\Delta t$] Passo de integração (0,5 dia no programa)
  \item[$z(t)$] Posição no plano da órbita escrita como número complexo, $z = x + \ii y$
\end{simbolos}
